\section{A correlation screen can discard a nonlinear cause}
\label{exec14:l4-nonlinear-screen}

The proposal should also evaluate what happens before the formal discovery
engine is called. A marginal correlation screen may reduce the number of
candidate pairs supplied to an agent, but its omissions constrain everything
downstream. The inspected Augur procedure uses a Spearman screen before
language judgments. Its reported forecasting performance does not establish
that this screen retains all intervention-relevant mechanisms
\citep{exec13_census_augur_arxiv}. An exact project example shows the
distinction without making any claim about an executed Augur benchmark.

Let $X\sim N(0,1)$ and $E\sim N(0,1)$ be independent, and set
\[
 Y=X^2-1+E.
\]
Then $E(Y\mid X=x)=x^2-1$, so $X$ and $Y$ are dependent, but
\[
 \operatorname{Cov}(X,Y)=E(X^3)-E(X)=0,
 \qquad \operatorname{Var}(Y)=E(X^4)-E(X^2)^2+1=3.
\]
The population Spearman correlation is also zero. To see this, the joint law
of $(X,Y)$ is unchanged when $X$ is replaced by $-X$. Since the normal
distribution function obeys $\Phi(-X)=1-\Phi(X)$, the expectation of
$(\Phi(X)-1/2)(F_Y(Y)-1/2)$ is equal to its negative and hence is zero.
The independent Gaussian error makes the outcome continuous, so the usual
continuous-rank interpretation applies. Neither a Pearson nor a Spearman
population threshold bounded away from zero retains this pair on the basis
of its marginal correlation.

The intervention response is not zero: replacing $X$ by $x$ gives
$E{Y\mid\operatorname{do}(X=x)\}=x^2-1$. Moving from zero to one changes
the mean by one; moving from minus one to one changes it by zero. Thus even
the phrase ``the effect of $X$'' is incomplete without the intervention
contrast. In a banking scenario, $X$ might be a standardized deviation from
a reference pricing point and the response might be curved. The example
does not establish that such a curve applies to deposits. It identifies a
specific failure regime that a discovery proposal should cover
\citep{exec14_discovery_screening_value}.

A deterministic feature $H=X^2-1$ would reveal the association in this law:
$\operatorname{Var}(H)=\operatorname{Cov}(H,Y)=2$. That feature is a useful
model representation, but it does not make $H$ an independently manipulable
business action. It maps $X$ and $-X$ to the same value. The interface must
retain the original action, the transformation and the intervention contrast.
Discovering relations among engineered features can change the graph's
variables and its feasible interventions. An agent should not relabel those
features as independently adjustable levers.

Several remedies have distinct costs. Retaining all candidate pairs avoids
this particular screen but increases search and testing burden. A nonlinear
dependence screen can detect other features of the joint law, subject to its
sampling assumptions and power. A scientifically motivated feature family
can make a restricted functional model plausible, but choosing features after
examining the final evaluation sample changes the procedure being assessed.
A language proposal of a nonlinear relation can be retained as a hypothesis
without being treated as verified evidence. The public prototype should
compare these alternatives with the same initial information and report
which candidate relations were removed before graph search began.

\section{Information gain and decision value can choose different queries}
\label{exec14:l4-information-objective}

Model disagreement or expected entropy reduction can be useful ways to select
a new experiment. The inspected Model Discovery Agent uses a disagreement
surrogate for information gain within its candidate mechanism pool
\citep{v12_la_2608_09696}. The banking decision adds another requirement:
the acquired information must concern something that can change a feasible
action. A high information score need not meet that requirement. The following
project example makes the difference exact.

Let a candidate world have two independent binary properties, $R$ and $S$.
Assume $\Pr(R=1)=1/10$ and $\Pr(S=1)=1/2$. Only $R$ affects the offer
decision: offering gives five units when $R=1$ and minus three when $R=0$;
not offering gives zero. Property $S$ concerns another graph feature that
does not affect this target. The current expected offer value is
$5/10-27/10=-11/5$, so the current action is no offer. Perfect observation
of $R$ would allow offering only in its favorable state, with expected value
$1/2$. Perfect observation of $S$ cannot change the action and has decision
value zero.

Nevertheless, a query revealing $S$ removes one bit of world uncertainty,
while a query revealing $R$ removes only
$h_2(1/10)\simeq .468996$ bits, where
$h_2(p)=-p\log_2p-(1-p)\log_2(1-p)$. Maximizing entropy reduction over
the full candidate world therefore prefers the decision-irrelevant query.
This is a difference between objectives, not a computational failure. It can
arise even when the candidate probabilities and the query likelihoods are
exactly known.

Now consider an informative but imperfect query $Z$ about $R$, with
$\Pr(Z=1\mid R=1)=4/5$ and $\Pr(Z=1\mid R=0)=1/5$. A positive answer
gives posterior probability $4/13$, and a negative answer gives $1/37$.
Both are below the offer threshold $3/8$, so the optimal action remains no
offer for either answer. This query has positive mutual information about
$R$, namely $h_2(13/50)-h_2(1/5)\simeq .104818$ bits, but zero value
for this decision. Information that changes beliefs can still be too weak to
cross a relevant action threshold.

A more accurate query has true-positive probability $19/20$ and
false-positive probability $1/20$. Its positive posterior is $19/28$ and
its negative posterior is $1/172$. Offering only after a positive answer
now gives ex ante value
\[
 5\left(\frac1{10}\frac{19}{20}\right)
 -3\left(\frac9{10}\frac1{20}\right)
 =\frac{17}{50}=.34.
\]
This remains below the perfect-information value $.5$. A query cost of $.1$
per comparable future decision would leave planning net value $.24$ under
these stipulations. The same priced irrelevant query has negative net value.
The probability and cost assumptions are authored inputs, not measured
calibration or vendor prices \citep{exec14_discovery_screening_value}.

The lesson for an agent is precise. Retain a model-discrimination objective
when scientific understanding is the declared target. For a specified banking
action, propagate candidate worlds into utilities and ask whether the query
can change the preferred action. If several bank decisions share the graph,
the utility may include several uses; their populations and weights need to
be stated rather than silently inferred. Customer constraints and acquisition
costs also belong in this calculation. Optimizing information gain is not
wrong, but its connection to investment value must be demonstrated.

The prototype can implement these finite reference problems before conducting
any expensive experiment. It should reproduce the different query rankings,
preserve the posterior update and explain the dependence on the prior and
loss. If a language-only query is assigned a likelihood, that likelihood
requires its own evidence; model confidence is not a substitute. A correct
implementation of this arithmetic would demonstrate decision accounting in
a declared finite problem. It would not establish that the candidate-world
set includes the bank's true mechanism or that a real proposed query has the
stipulated reliability.
\label{exec14:l4-final-end}
